% Lean module: MonadFormal.Temporal.Variational.Shift
\subsubsection{The shifted-difference adjoint}
\label{sec:variational-shift}

For $0\le\delta\le1$, define the nonlocal force
\begin{equation}\label{eq:variational-force}
 G_\delta[x](t)
 =\mathbf1_{[\delta,1]}(t)\bigl(x(t)-x(t-\delta)\bigr)
  -\mathbf1_{[0,1-\delta]}(t)\bigl(x(t+\delta)-x(t)\bigr).
\end{equation}
The indicators are part of the operator. They exclude a shifted
neighbour that lies outside the interval; no periodic or exterior
boundary value is supplied.
\leandef{MonadFormal.Temporal.Variational.Interval.shiftForce}

\begin{lemma}[Shifted pairing with its boundary supports]
\label[lemma]{lem:shift-pairing}
For continuous real functions $x,h$,
$G_\delta[x]\in L^2(\mu)$ and
\[
 \int G_\delta[x](t)h(t)\,d\mu(t)
 =\int\bigl(x(t+\delta)-x(t)\bigr)
          \bigl(h(t+\delta)-h(t)\bigr)\,d\nu_\delta(t).
\]
\end{lemma}
\begin{proof}
Continuous functions are bounded on $[0,1]$, so the two restricted
differences are square-integrable. Split the product on the right.
In the term containing $h(t+\delta)$, set $s=t+\delta$:
\[
 \int_0^{1-\delta}\!\bigl(x(t+\delta)-x(t)\bigr)h(t+\delta)\,dt
 =\int_\delta^1\!\bigl(x(s)-x(s-\delta)\bigr)h(s)\,ds.
\]
The other term remains on $[0,1-\delta]$. Writing both as
integrals on $[0,1]$ introduces precisely the two indicators in
\eqref{eq:variational-force}. Endpoints have Lebesgue measure zero.
\leanproof{MonadFormal.Temporal.Variational.Interval.continuous_memLp_unit,MonadFormal.Temporal.Variational.Interval.shiftForce_memLp,MonadFormal.Temporal.Variational.Interval.shifted_pairing,MonadFormal.Temporal.Variational.Interval.unit_integral_eq_interval,MonadFormal.Temporal.Variational.Interval.unit_integral_indicator,MonadFormal.Temporal.Variational.Interval.shiftForce_pairing}
\end{proof}
